% ── Imigrace, agenturní zaměstnávání a platformová práce ─────────────────────

Pracovní podmínky a~mzdovou distribuci v~\acgen{ČR} ovlivňuje také imigrace a~nestandardní
formy zaměstnávání. Kombinace tří souběžných tlaků ---~migrační vlny z~Ukrajiny, inklinace
v~populaci k~vysokému podílu \acs{OSVČ} často v~režimu zastřené závislé práce (\emph{švarcsystému})
a~nově nastupující segment platformové práce, vytváří paralelní trh
práce, na~který se~tradiční mechanismy \acgen{KV} dosud uplatňují jen
omezeně. Pro~sociální dialog v~\acgen{ČR}, již dnes oslabený nízkým pokrytím \acs{KS},
představuje tato situace riziko eroze dosahu stávajících \aclpgen{KS}
u~rostoucích segmentů pracovního trhu, kterých se~vyjednané podmínky netýkají.\par

Polents a~Dvořáková v~empirické studii dokumentují systematickou degradaci kvalifikace
ukrajinských uprchlíků s~dočasnou ochranou na~českém trhu práce: navzdory tomu,
že~přibližně \SI{70}{\percent} respondentek má~vysokoškolské vzdělání, pouze
\SI{7,1}{\percent} pracuje na~pozici odpovídající kvalifikaci; \SI{17,9}{\percent}
v~plnoúvazkovém zaměstnání pod úrovní kvalifikace, \SI{25}{\percent} v~šedé
ekonomice ("na ruku") a~plných \SI{50}{\percent} bez zaměstnání. Hlavními
bariérami jsou jazyk, nedostatečně rozvinutá nostrifikace zahraničních diplomů
a~nedostatek kapacit předškolní péče (populace migrantů
tvoří z~\SI{99,1}{\percent} ženy, často s~nezletilými dětmi). Z~makroekonomického
pohledu představuje ukrajinská diaspora podstatný stabilizátor demografické
bilance ---~bez ní by~\acgen{ČR} v~roce~2023 vykazovala čistý úbytek
populace, ze~strukturálního pohledu však koncentrace migrantů
do~nízkokvalifikovaných sektorů (gastronomie, úklid, sklady) generuje
dumpingový mzdový tlak a~zároveň posiluje segment, na~který \ac{KV} má dosud
minimální vliv. Autorky studie explicitně volají po~zapojení sociálních partnerů
do~tvorby integrační politiky jako prevenci oddělování „českých” a~„ukrajinských” segmentů trhu
práce.~\cite{PolentsDvorakova_UkraineRefugees2024}\par

Migrační a~subdodavatelský rozměr nelze oddělit od~evropského režimu volného pohybu. Schengenský prostor a~navazující pravidla koordinace sociálního zabezpečení snižují transakční náklady mobility pro~zaměstnance, firmy i~přeshraniční zadavatele, což je~pro \acacc{ČR} jako otevřenou ekonomiku dvojsečné. Na~jedné straně umožňují českým pracovníkům využívat vyšší mzdovou hladinu v~sousedních regionech \aclgen{DE} a~\aclgen{AT}; na~straně druhé zvyšují tlak na~domácí zaměstnavatele v~pohraničí a~rozšiřují prostor pro~agenturní práci a~subdodavatelské řetězce. Bez~domácí odvětvové mzdové koordinace se~tak mobilita nemusí proměnit v~konvergenci pracovních standardů, ale v~prohloubení mzdové opony mezi jádrem a~periferií \acgen{EU}.~\cite{EC_WhitePaper_FutureOfEurope_2017}\par

Druhou velkou výjimkou z~působnosti \acgen{KV} jsou \acs{OSVČ}
v~ekonomicky závislém postavení. Eurostat řadí \acgen{ČR} mezi ekonomiky s~nejvyšším podílem \acs{OSVČ} v~\acs{EU}: podíl samostatně výdělečně činných
osob na~zaměstnanosti se~dlouhodobě drží kolem \SI{17}{\percent} oproti
průměru \acs{geo-EU27} přibližně \SI{14}{\percent}. Významnou součástí tohoto
podílu je~\emph{švarcsystém} ---~výkon práce splňující všechny znaky závislé
činnosti dle §~2~zákoníku práce, ovšem na~bázi obchodněprávní smlouvy
s~\acs{OSVČ}. Odhaduje se, že počet takto zaměstnaných přesahuje 100~tisíc (cca 3~\%) a~roste. Hlavní motivací je~asymetrický daňový klín: zatímco zaměstnanec
inkasuje méně než \SI{60}{\percent} superhrubé mzdy, \acs{OSVČ} s~typickým 60\,\%
získá řádově \SI{15}{\percent} při minimálních odvodech a~zdanění příjmu
(\SI{100}{\czk} při mediánu mezd). Paušální daňový režim zavedený od~roku~2021 snižuje administrativu a~zvyšuje krátkodobý čistý příjem části \acs{OSVČ}, avšak zároveň může vést k~nižší účasti na~nemocenském pojištění, které je pro \acs{OSVČ} narozdíl od zaměstnanců dobrovolné. Například peněžitá pomoc v~mateřství patří do~systému nemocenského pojištění; faktická možnost čerpat ochranu v~období rodičovství proto závisí na~předchozí dobrovolné účasti a~výši plateb. V~odvětvích, kde je samostatná výdělečná činnost využívána jako náhrada pracovního poměru, se~tak daňová arbitráž neprojevuje jen v~důchodech, ale i~v~nerovném rozložení rizik péče, nemoci a~mateřství.
Z~hlediska sociálního dialogu má~tento systém tři důsledky: erozi systému sociálního pojištění a~tedy
i~financovatelnosti budoucích důchodů; vyloučení velké části pracovních
vztahů z~oblasti, na~kterou se~vztahuje právo \acgen{KV}; a~asymetrické
postavení \acs{OSVČ} při sjednávání podmínek.
Pokyny Komise z~roku~2022 tuto asymetrii částečně narovnávají tím, že~připouští
\acgen{KV} pro~ekonomicky závislé \acs{OSVČ} za určitých podmínek.~\cite{EC_Guidelines_CB_SoloSelfEmployed_2022}~\cite{CSSZ_RocenkaOSVC_2024}~\cite{PAQ_Svarcsystem}\par

Třetím rostoucím segmentem je~platformová práce. Podle podkladových
studií \acs{MPSV} pracuje v~\acgen{ČR} přes digitální platformy přibližně
145~tis.\ osob měsíčně (kurýři, taxi, řemeslné a~osobní služby
zprostředkované online). Návrh zákona o~platformové práci (\emph{ZOPP})
představený \acs{MPSV} v~březnu~2026 transponuje směrnici \acs{EU}~2024/2831
o~zlepšení pracovních podmínek při práci prostřednictvím
platforem.
Návrh současně zasahuje do~obecné definice závislé práce: dosavadní kumulativní test se~má soustředit na~vztah nadřízenosti a~podřízenosti, práci jménem zaměstnavatele a~konkrétní projevy organizace, dohledu, pokynů a~pracovní doby; osobní výkon práce se~přesouvá z~definičního znaku do~podmínky výkonu.
Klíčovým právním nástrojem je~vyvratitelná právní doměnka
zaměstnaneckého poměru: pokud existují skutečnosti důvodně značící naplnění
znaků závislé práce, je~na~platformě, aby~prokázala opak. Důležitým
interpretačním vodítkem pro~rozlišení autonomie versus závislosti je~rozhodnutí
\acf{SDEU}~C-692/19 \emph{Yodel Delivery}, které definovalo flexibilitu pracovní
doby jako~podstatný indikátor skutečné autonomie pracovníka.
Paralelně směrnice~(\acs{EU})~2018/957 o~vysílání pracovníků zpřesňuje
podmínky přeshraničního dočasného nasazení v~rámci \acs{EU} ---~oblast,
která se~na~\acgen{ČR} dotýká primárně příjmově (české firmy vysílající
zaměstnance do~Německa a~Rakouska). V~subdodavatelské ekonomice je~tato úprava podstatná proto, že mobilita zakázky, pracovníků a~odpovědnosti se~snadno odděluje: formální zaměstnavatel může být český subdodavatel, faktická ekonomická kontrola leží u~zahraničního odběratele a~pracovní standard se~pohybuje mezi českou a~cílovou úpravou. Bez~silných \acs{KSVS} se~tak přeshraniční mobilita může stát kanálem mzdové arbitráže namísto konvergence pracovních podmínek.~\cite{MPSV_ZakonOdmenovani_2026}~\cite{Kabat_ZOPP_2026}~\cite{smernice_posted_workers_2018_957}\par

Souběh všech tří faktorů ---~migrace, švarcsystému a~platformové práce
---~tvoří výzvu pro~budoucnost \acgen{KV} v~\acgen{ČR}. Pokud se působnost
\acp{KSVS} a~\acp{KS} nerozšíří nad rámec klasického pracovního poměru,
budou se~dlouhodobě vztahovat ke~zmenšujícímu se~jádru pracovní síly,
zatímco rostoucí segmenty alternativních forem práce zůstanou bez kolektivní ochrany
a~bez příspěvku k~financování systémových mechanismů (sociální pojištění, \acgen{APZ},
sektorové fondy vzdělávání). Reformní agenda by~proto měla zahrnovat:
\begin{itemize}
  \item rozšíření aplikační působnosti \acs{KSVS} na~ekonomicky závislé \acs{OSVČ} v~odvětvích s~prokazatelně systematickým výskytem švarcsystému;
  \item explicitní zařazení platformových pracovníků do~působnosti odvětvových \acs{KSVS} prostřednictvím definic v~\emph{ZOPP};
  \item strukturované zapojení sociálních partnerů do~tvorby integrační politiky imigrantů.
\end{itemize}
Bez těchto adaptací riskuje sociální dialog ztrátu reprezentativity vůči rostoucí
části pracovní síly.\par
